\subsection{群的表现}

设 G 为一个群，且 \(\mathbf{t} = (t_{i})_{i \in I}\) 是 G 的一个元素族。设 \(f_{t}\) 为自由群 F(I) 到 G 的、把 i 映到 \(t_{i}\) 的唯一同态。 \(f_{t}\) 的像是由元素 \(t_{i}\) 生成的 G 的子群。 \(f_{t}\) 的核的元素称为族 t 的关系子。若 \(f_{t}\) 是满射（相应地单射、双射），则称 t 为生成族（相应地自由族、基本族）。

设 G 为一个群。G 的一个表现是一个有序对 \((\mathbf{t}, \mathbf{r})\) ，它由一个生成族 \(\mathbf{t} = (t_{i})_{i \in I}\) 和一个关系子族 \(\mathbf{r} = (r_{j})_{j \in J}\) 组成，使得 \(N_{t}\) （即 \(f_{t}\) 的核）由元素 \(gr_{j}g^{-1}\) （对 \(g \in \mathrm{F}(\mathrm{I})\) 和 \(j \in J\) ）生成。这等价于说 \(N_{t}\) 是由 \(r_{j}\) （对 \(j \in J\) ）生成的 F(I) 的正规子群（换言之，是 F(I) 中包含元素 \(r_{j}\) ( \(j \in J\) ) 的最小正规子群；参见 § 4，第 4 号）。以语言上的滥用，称生成元 \(t_{i}\) 和关系 \(r_{j}(\mathbf{t}) = e\) 构成群 G 的一个表现。

设 I 为一个集合，且 \(\mathbf{r} = (r_j)_{j \ensuremath{\in} J}\) 为自由群 F(I) 的一族元素。设 N(r) 为由 \(r_j\) （对 \(j \ensuremath{\in} J\) ）生成的 F(I) 的正规子群。设 F(I, r) = F(I)/N(r)，并记 \(\tau_i\) 为 \(i\) 模 N(r) 的类。有序对 \((\tau, \mathbf{r})\) （其中 \(\tau = (\tau_i)_{i \ensuremath{\in} I}\) ）是群 F(I, r) 的一个表现；若 G 是一个群且 (t, r) 是 G 的一个表现，其中 t = (t\_i)\_\{i \ensuremath{\in} I\}，则存在唯一的同构 \(u\) ，它把 F(I, r) 映到 G 上，使得 \(u(\tau_i) = t_i\) 对所有 \(i \ensuremath{\in} I\) 成立。称群 F(I, r) 由生成元 \(\tau_i\) 和关系子 \(r_j\) 定义，或者以语言上的滥用，称它由生成元 \(\tau_i\) 和关系 \(r_j(\tau) = e\) 定义。当 I = {[}1, n{]} 且 J = {[}1, m{]}时，称 F(I, r) 由该表现定义

\[
\langle \tau_ {1}, \dots , \tau_ {n}; r _ {1}, \dots , r _ {m} \rangle .
\]

若 \(r_j = u_j^{-1}v_j\) ，其中 \(u_j\) 和 \(v_j\) 属于 \(\mathrm{F(I)}\) ，则该表现同样用符号

\[
\langle \tau_ {1}, \dots , \tau_ {n}; u _ {1} = v _ {1}, \dots , u _ {m} = v _ {m} \rangle .
\]

例子。(1) 由表现 \(\langle \tau; \tau^q = e \rangle\) 定义的群是 \(q\) 阶循环群。

\begin{enumerate}
\def\labelenumi{(\arabic{enumi})}
\setcounter{enumi}{1}
\tightlist
\item
  由表现 \(\langle x, y; xy = yx \rangle\) 定义的群同构于 \(\mathbf{Z} \times \mathbf{Z}\) 。
\end{enumerate}

命题 9. 设 G 是一个群， \(\mathbf{t} = (t_i)_{i \in I}\) 是 G 的一个生成族， \(\mathbf{r} = (r_j)_{j \in J}\) 是 \(\mathbf{t}\) 的一个关系子族。以下条件等价：

\begin{enumerate}
\def\labelenumi{(\alph{enumi})}
\item
  有序对 \((\mathbf{t},\mathbf{r})\) 是 G 的一个表现。
\item
  设 \(\mathbf{G}'\) 为一个群，且 \(\mathbf{t}' = (t_i')_{i \in \mathbf{I}}\) 为 \(\mathbf{G}'\) 的一个元素族。若 \(r_j(\mathbf{t}') = e\) 对所有 \(j \in \mathbf{J}\) 成立，则存在一个同态 \(u\) ，它把 \(\mathbf{G}\) 映入 \(\mathbf{G}'\) ，使得 \(u(t_i) = t_i'\) 对所有 \(i \in \mathbf{I}\) 成立。
\item
  设 \(\overline{\mathbf{G}}\) 为一个群，且 \(\overline{\mathbf{t}} = (\bar{t}_i)_{i\in \mathbf{I}}\) 为 \(\overline{\mathbf{G}}\) 的一个生成族，满足 \(r_j(\overline{\mathbf{t}}) = e\) 对所有 \(j\in \mathbf{J}\) 成立。每一个把 \(\overline{\mathbf{G}}\) 映入 \(\mathbf{G}\) 、且把 \(\bar{t}_i\) 映到 \(t_i\) （对所有 \(i\in \mathbf{I}\) ）的同态都是同构。
\end{enumerate}

设 \(f\) 表示 \(\mathbf{F}(\mathbf{I})\) 到 \(\mathbf{G}\) 中的、把 \(i\) 映到 \(t_i\) （对所有 \(i \in \mathbf{I}\) ）的同态，并设 \(\mathbf{N}\) 为 \(f\) 的核。

\begin{enumerate}
\def\labelenumi{(\alph{enumi})}
\item
  \(\Rightarrow\) (b)：假设 \((\mathbf{t},\mathbf{r})\) 是 G 的一个表现，且令 \(\mathbf{t}' = (t_i')_{i\in I}\) 是某个群 \(\mathbf{G}'\) 的元素的一个族，满足 \(r_j(\mathbf{t}') = e\) 对所有 \(j\in J\) 成立。设 \(f'\) 是 F(I) 到 \(\mathbf{G}'\) 中的、以 \(f'(i) = t_i'\) （对所有 \(i\in I\) ）为特征的同态。根据假设， \(f'(r_j) = e\) 对所有 \(j\in J\) 成立，且由于 N 由元素 \(gr_jg^{-1}\) （对 \(j\in J\) 和 \(g\in \mathrm{F}(\mathrm{I}),f'(\mathrm{N}) = \{e\}\) ）生成。由于同态 \(f\colon \mathrm{F}(\mathrm{I})\to \mathrm{G}\) 是满射且核为 N，因此存在一个同态 \(u\colon \mathrm{G}\to \mathrm{G}'\) ，使得 \(f' = u\circ f\) 。于是 \(u(t_i) = u(f(i)) = f'(i) = t_i'\) 。
\item
  \(\Rightarrow\) (c)：假设条件 (b) 成立。设 \(\mathbf{t} = (t_i)_{i\in I}\) 是群 G 的一个生成族，满足 \(r_j(\mathbf{t}) = e\) 对所有 \(j\in J\) 成立，并设 \(v\) 是 \(\overline{\mathbf{G}}\) 到 G 中的同态，满足 \(v(t_i) = t_i\) 对所有 \(i\in I\) 成立。由于族 \((t_i)_{i\in I}\) 生成 G，同态 \(v\) 是满射的。根据性质 (b)，存在一个同态 \(u\colon G\to \overline{\mathbf{G}}\) ，使得 \(u(t_i) = \bar{t}_i\) 对所有 \(i\in I\) 成立。则 \(u(v(\bar{t}_i)) = \bar{t}_i\) 对所有 \(i\in I\) 成立，因此 \(u \circ v\) 是 \(\overline{G}\) 上的恒等映射，这证明了 \(v\) 是单射。于是 \(v\) 是一个同构，且条件 (c) 成立。
\item
  \(\Rightarrow\) (a)：假设条件 (c) 成立。设 \(t_i'\) 为 \(i\) 在 \(\mathbf{F}(\mathbf{I},\mathbf{r})\) 中的典范像，并设 \(\mathbf{t}' = (t_i')_{i\in \mathbf{I}}\) ；于是 \(r_j(\mathbf{t}') = e\) 对所有 \(j\in \mathbf{J}\) 成立。由于 \(r_j(\mathbf{t}) = e\) 对所有 \(j\in \mathbf{J}\) 成立，存在且仅存在一个同态 \(v\) ，它把 \(\mathbf{F}(\mathbf{I},\mathbf{r})\) 映入 \(\mathbf{G}\) ，使得 \(v(\bar{t}_i) = t_i\) 对所有 \(i\in \mathbf{I}\) 成立。根据 (c)， \(v\) 是 \(\mathbf{F}(\mathbf{I},\mathbf{r})\) 到 \(\mathbf{G}\) 上的一个同构，它把表现 \((\mathbf{t}',\mathbf{r})\) （\(\mathbf{F}(\mathbf{I},\mathbf{r})\) 的）变为表现 \((\mathbf{t},\mathbf{r})\) （\(\mathbf{G}\) 的）。
\end{enumerate}

